\chapter{صيغة التغاير التكرارية}

\section{مقدمة}

تعتبر هذه الطريقة التي اقترحها العالم He في اواخر التسعينات من اقوى الادوات لحل المعادلات التفاضلية الخطية وغير الخطية.

\section{الفكرة الاساسية لهذه الطريقة}
تعتمد هذه الطريقة على بناء مؤثر تصحيحي 
(functional correction)
يقوم بتقريب الحل تدريجياً من خلال عملية تكرارية دون الحاجة لتحويل "لابلاس" او "متسلسلات تايلور" المعقدة في البداية.

\section{الصيغة الرياضية للطريقة}
لنفرض وجود معادلة تفاضلية عامة بالصيغة
\[
L\{u(t)\} + N\{u(t)\} = g(t)
\]
حيث
\begin{itemize}[nosep]
	\item $L$ المؤثر الخطي (linear operator).
	\item $N$ المؤثر غير الخطي (non-linear operator).
	\item $g(t)$ الدالة المصدرية (source term).
\end{itemize}

\subsection*{بناء المؤثر التصحيحي}
يتم بناء التكرار وفق الصيغة التالية
\[
u_{n+1}(t) = u_n(t) + \int_{0}^{t} \lambda(\tau) \Big[L u_n(\tau) + N \tilde{u}_n(\tau) - g(\tau)\Big] d\tau
\]
حيث
\begin{itemize}[nosep]
	\item $\lambda$ : هو مضروب لاكرانج العام 
	(general lagrange multiplier) ويتم ايجاده بأستخدام نظرية التغاير.
	
	\item $\tilde{u}_n$ : يمثل التغاير المقيد 
	(restricted variation) ، اي يتم اعتباره ثابتاً عند اشتاقاق $\lambda$.
\end{itemize}

\subsection*{خطوات الحل}
\begin{enumerate}[label=\textbf{\arabic*.}]
	\item تحديد مضروب لاكرانج $\lambda$ الامثل للمعادلة المعطاة.
	\item تعويض قيمة $\lambda$ في الصيغة التكرارية.
	\item اختيار دالة ابتدائية $u_0(t)$ (غالباً ما تكون الشروط الابتدائية).
	\item اجراء التكرارات ($n=0,1,2,\dots$) حتى الوصول الى الحل التقريبي.
\end{enumerate}

\begin{example}
	لنأخذ المعادلة التفاضلية البسيطة
	\[
	u'(t) = t, \quad u(0) = 0
	\]
	\textbf{اولاً:} نحدد العناصر 
	\begin{itemize}
		\item المؤثر الخطي $L(u) = u'(t)$.
		\item لا يوجد مؤثر غير خطي ، اي ان $N=0$.
		\item دالة المصدر $g(t) = t$
	\end{itemize}
	\textbf{ثانياً:} صياغة دالة التصحيح
	\[
	u_{n+1}(t) = u_n(t) + \int_{0}^{t} \lambda(\tau) \Big[L u_n(\tau) + N \tilde{u}_n(\tau) - g(\tau)\Big] d\tau
	\]
	\[
	\Rightarrow
	u_{n+1}(t) = u_n(t) + \int_{0}^{t} \lambda(\tau) \Big[u_n'(\tau)   - \tau\Big] d\tau
	\]
	\textbf{ثالثاً:} اختيار مضاعف لاكرانج ، حيث في هذه الحالة البسيطة المضاعف المناسب هو $\lambda(\tau) = -1$.
	اذن العلاقة التكرارية تصبح
	\[
	u_{n+1}(t) = u_n(t) - \int_{0}^{t}  \Big[u_n'(\tau)   - \tau\Big] d\tau
	\]
	\textbf{رابعاً:} البدء بالتقريب الاولي ، نبدأ بالحل الابتدائي 
	$u_0(t) = 0$\\
	\textbf{خامساً:} نجد التكرار الأول
	\begin{align*}
		u_{1}(t) &= u_0(t) - \int_{0}^{t} \Big[u_0'(\tau)   - \tau\Big] d\tau\\
		& = 0 - \int_{0}^{t} 0 - \tau d\tau\\
		& = \int_{0}^{t} \tau d\tau \\
		& = \frac{t^2}{2}
	\end{align*}
	الآن نجد التكرار الثاني
	\begin{align*}
		u_{2}(t) &= u_1(t) - \int_{0}^{t} \Big[u_1'(\tau)   - \tau\Big] d\tau\\
		& = \frac{t^2}{2} - \underbrace{\int_{0}^{t} (\tau - \tau) d \tau}_{ = 0} = \frac{t^2}{2}
	\end{align*}
	نلاحظ ان 
	$u_1(t) = u_2(t)$ اذن يمكننا الاستنتاج ان
	\[
	u(t) = u_1(t) = u_2(t) = \cdots = \frac{t^2}{2}
	\]
\end{example}

\begin{example}
	لنأخذ المعادلة التفاضلية غير الخطية
	\[
	\frac{dy}{dx} + y^2 = 1, \quad y(0)=0
	\]
	\textbf{اولاً:} نحدد العناصر 
	\begin{itemize}
		\item المؤثر الخطي $L(y) = \dfrac{dy}{dx}$.
		\item المؤثر غير الخطي $N(y) = y^2$.
		\item دالة المصدر $g(x) = 1$.
	\end{itemize}
	\textbf{ثانياً:} صياغة دالة التصحيح
	\[
	y_{n+1}(x) = y_n(x) + \int_{0}^{x} 
	\lambda(\tau) 
	\Big[
	L y_n(\tau) + N \tilde{y}_n(\tau) - g(\tau)
	\Big] d\tau
	\]
	\[
	\Rightarrow
	y_{n+1}(x) = y_n(x) + \int_{0}^{x} 
	\lambda(\tau) 
	\Big[
	y_n'(\tau) + y_n^2(\tau) - 1
	\Big] d\tau
	\]
	\textbf{ثالثاً:} اختيار مضاعف لاكرانج.  
	بإجراء التباين نحصل على المضاعف المناسب
	\[
	\lambda(\tau) = -1
	\]
	اذن العلاقة التكرارية تصبح
	\[
	y_{n+1}(x) =
	y_n(x) -
	\int_{0}^{x}
	\Big[
	y_n'(\tau) + y_n^2(\tau) - 1
	\Big] d\tau
	\]
	\textbf{رابعاً:} نبدأ بالتقريب الابتدائي
	\[
	y_0(x)=0
	\]
	\textbf{خامساً:} التكرار الأول
	\begin{align*}
		y_1(x)
		&= y_0(x) -
		\int_{0}^{x}
		\Big[
		y_0'(\tau) + y_0^2(\tau) - 1
		\Big] d\tau \\
		&= 0 - \int_{0}^{x} (0 + 0 - 1)\, d\tau \\
		&= \int_{0}^{x} 1\, d\tau \\
		&= x
	\end{align*}
	\textbf{التكرار الثاني}
	\begin{align*}
		y_2(x)
		&= y_1(x) -
		\int_{0}^{x}
		\Big[
		y_1'(\tau) + y_1^2(\tau) - 1
		\Big] d\tau \\
		&= x -
		\int_{0}^{x}
		\Big[
		1 + \tau^2 - 1
		\Big] d\tau \\
		&= x - \int_{0}^{x} \tau^2 d\tau \\
		&= x - \frac{x^3}{3}
	\end{align*}
\textbf{التكرار الثالث}
\begin{align*}
	y_3(x)
	&= y_2(x) -
	\int_{0}^{x}
	\Big[
	y_2'(\tau) + y_2^2(\tau) - 1
	\Big] d\tau
\end{align*}
نحسب أولاً المشتقة والمربع:
\[
y_2(x)=x-\frac{x^3}{3}
\]
\[
y_2'(\tau)=1-\tau^2
\]
\[
y_2^2(\tau)=\left(\tau-\frac{\tau^3}{3}\right)^2
= \tau^2-\frac{2}{3}\tau^4+\frac{1}{9}\tau^6
\]
اذن 
\begin{align*}
	y_2'(\tau)+y_2^2(\tau)-1
	&=(1-\tau^2)+\left(\tau^2-\frac{2}{3}\tau^4+\frac{1}{9}\tau^6\right)-1 \\
	&= -\frac{2}{3}\tau^4+\frac{1}{9}\tau^6
\end{align*}
إذن
\begin{align*}
	y_3(x)
	&= x-\frac{x^3}{3}
	- \int_{0}^{x}
	\left(
	-\frac{2}{3}\tau^4+\frac{1}{9}\tau^6
	\right)d\tau \\
	&= x-\frac{x^3}{3}
	+\int_{0}^{x}
	\left(
	\frac{2}{3}\tau^4-\frac{1}{9}\tau^6
	\right)d\tau
\end{align*}
نحسب التكامل:
\[
\int_{0}^{x}\frac{2}{3}\tau^4 d\tau=\frac{2}{15}x^5
\]
\[
\int_{0}^{x}\frac{1}{9}\tau^6 d\tau=\frac{1}{63}x^7
\]
إذن الحل الثالث:
\[
y_3(x)
=
x-\frac{x^3}{3}
+\frac{2}{15}x^5
-\frac{1}{63}x^7
\]
	وبالاستمرار نحصل على متسلسلة تقارب الحل الدقيق
	\[
	y(x) = \tanh(x)
	\]
\end{example}

\begin{example}
	لنأخذ المعادلة التفاضلية الجزئية غير الخطية
	\[
	u_t - 6u u_x + u_{xxx} = 0, \quad u(x,0)=6x
	\]
	\textbf{اولاً:} نحدد العناصر
	\begin{itemize}
		\item المؤثر الخطي $L(u)=u_t+u_{xxx}$
		\item المؤثر غير الخطي $N(u)=-6u u_x$
		\item لا يوجد حد مصدر، أي $g=0$
	\end{itemize}
	\textbf{ثانياً:} صياغة دالة التصحيح
	\[
	u_{i+1}(x,t)=u_i(x,t)+\int_{0}^{t}\lambda(\tau)\Big[(u_i)_{\tau}+(u_i)_{xxx}-6u_i (u_i)_x\Big]d\tau
	\]
	\textbf{ثالثاً:} باجراء التباين نحصل على مضاعف لاكرانج
	\[
	\lambda(\tau)=-1
	\]
	اذن العلاقة التكرارية تصبح
	\[
	u_{i+1}(x,t)=u_i(x,t)-\int_{0}^{t}\Big[(u_i)_{\tau}+(u_i)_{xxx}-6u_i (u_i)_x\Big]d\tau
	\]
	\textbf{رابعاً:} التقريب الابتدائي من الشرط الابتدائي
	\[
	u_0(x,t)=6x
	\]
	\textbf{خامساً:} التكرار الأول
	\[
	(u_0)_{\tau}=0,\quad (u_0)_x=6,\quad (u_0)_{xxx}=0
	\]
	\[
	-6u_0 (u_0)_x=-6(6x)(6)=-216x
	\]
	\[
	u_1(x,t)=6x-\int_{0}^{t}(-216x)d\tau
	\]
	\[
	u_1(x,t)=6x+216xt
	\]
	\textbf{التكرار الثاني}
	\[
	u_2(x,t)=u_1(x,t)-\int_{0}^{t}\Big[(u_1)_{\tau}+(u_1)_{xxx}-6u_1 (u_1)_x\Big]d\tau
	\]
	حيث
	\[
	u_1(x,t)=6x+216xt
	\]
	نحصل بعد التعويض والتكامل على
	\[
	u_2(x,t)=6x+216xt+3888xt^2+279936xt^3
	\]
	\textbf{التكرار الثالث}
	\[
	u_3(x,t)=u_2(x,t)-\int_{0}^{t}\Big[(u_2)_{\tau}+(u_2)_{xxx}-6u_2 (u_2)_x\Big]d\tau
	\]
	وبعد الحساب نحصل على
	\[
	u_3(x,t)=6x+216xt+3888xt^2+279936xt^3+10077696xt^4+2176782336xt^5
	\]
\end{example}